\cotefile{83}
\scanpage{32}\mspage{2}

(10) de $H$. On pose également

\[
L_0^{\deleted{\unsure{\prime|1}}}={L'_1}^{-1}L'_0,\qquad
L_{\unsure{1|3}}^{\deleted{\unsure{\prime|1}}}={L'_0}^{-1}L'_1,\tag*{(11)}
\]

ce sont donc des modules inversibles, canoniquement inverses l'un de l'autre. On a des homomorphismes canoniques

\[
\begin{aligned}
 e_0:L_0&\longrightarrow H,&e_1:L_1&\longrightarrow H\\
 \beta&\longmapsto \bigl(\begin{smallmatrix}1&\beta\\0&1\end{smallmatrix}\bigr),& \gamma&\longmapsto \bigl(\begin{smallmatrix}1&0\\\gamma&1\end{smallmatrix}\bigr)
\end{aligned}\tag*{(12)}
\]

et on pose

\[
\begin{aligned}
 U_0&=e_0(L_0)=\left\{\bigl(\begin{smallmatrix}1&\beta\\0&1\end{smallmatrix}\bigr)\mid\beta\in L_0\right\}\\
 U_1&=e_1(L_1)=\left\{\bigl(\begin{smallmatrix}1&0\\\gamma&1\end{smallmatrix}\bigr)\mid\gamma\in L_1\right\}.
\end{aligned}\tag*{(13)}
\]

On a pour $\lambda\in\Lambda^*$, $g=\bigl(\begin{smallmatrix}a&b\\c&d\end{smallmatrix}\bigr)\in\deleted{\unsure{\mathrm{SL}|H}}G$,

\[
\left\{\begin{aligned}
 u(\lambda)g&=\bigl(\begin{smallmatrix}\lambda a&\lambda b\\\lambda^{-1}c&\lambda^{-1}d\end{smallmatrix}\bigr)&
 gu(\lambda)&=\bigl(\begin{smallmatrix}\lambda a&\lambda^{-1}b\\\lambda c&\lambda^{-1}d\end{smallmatrix}\bigr)\\
 u(\lambda)gu(\lambda)^{-1}&=\bigl(\begin{smallmatrix}a&\lambda^2b\\\lambda^{-2}c&d\end{smallmatrix}\bigr)
\end{aligned}\right.\tag*{(14)}
\]

en particulier

\[
\left\{\begin{aligned}
 u(\lambda)(e_0(\beta))&=e_0(\lambda^2\beta)\\
 u(\lambda)(e_1(\gamma))&=e_0(\lambda^{-2}\gamma)
\end{aligned}\right.\tag*{(15)}
\]

\margin{attention : l'exposant $2$ des $\lambda$ dans le $2^{\mathrm e}$ membre !}

Ainsi $T$ normalise $U_0$ et $U_1$, et on a\deleted{\unsure{t}} bien sûr

\[
T\cap U_0=T\cap U_1=\{1\},\tag*{(16)}
\]

on introduit

\[
\left\{\begin{aligned}
 B_0&=T.U_0=\{e_0(\beta)u(\lambda)\mid\beta\in\Lambda,\lambda\in\Lambda^*\}
 =\left\{\bigl(\begin{smallmatrix}\lambda&\lambda^{-1}\beta\\0&\lambda^{-1}\end{smallmatrix}\bigr)\;\middle|\;\begin{smallmatrix}\lambda\in\Lambda^*\\\beta\in\Lambda\end{smallmatrix}\right\}\\
 &\inserted{=\left\{\bigl(\begin{smallmatrix}a&b\\c&d\end{smallmatrix}\bigr)\in H\mid c=0\right\}}\\
 &\deleted{\unsure{B_1}}\\
 B_1&=T.U_1=\{u(\lambda)e_1(\gamma)\mid\gamma\in\Lambda,\lambda\in\Lambda^*\}
 =\left\{\bigl(\begin{smallmatrix}\lambda&0\\\lambda^{\unsure{-|-1}}\gamma&\lambda^{-1}\end{smallmatrix}\bigr)\;\middle|\;\begin{smallmatrix}\lambda\in\Lambda^*\\\gamma\in\Lambda\end{smallmatrix}\right\}\\
 &\ =\left\{\bigl(\begin{smallmatrix}a&b\\c&d\end{smallmatrix}\bigr)\in H\mid b=0\right\}.
\end{aligned}\right.\tag*{(17)}
\]

qui sont des sous-groupes de $H$ normalisant $U_0$ resp. $U_1$. On a

\[
B_0\cap B_1=T=\left\{\bigl(\begin{smallmatrix}a&b\\c&d\end{smallmatrix}\bigr)\in H\mid b=0\deleted{\unsure{=}}c=0\right\}.\tag*{(18)}
\]

% archive sheet number: 31
\scanpage{33}

On introduit aussi une application (\uline{pas} un homomorphisme !)

\[
\left\{\begin{aligned}
 \deleted{\unsure{v}}\sigma_0 &:L_0^*\longrightarrow H\\
 \sigma_0(\beta)&\deleted{\longmapsto}=\bigl(\begin{smallmatrix}0&\deleted{\unsure{-}}\beta^{\deleted{\unsure{-1}}}\\-\beta^{-1}&0\end{smallmatrix}\bigr)
\end{aligned}\right.\tag*{(19)}
\]

\margin{\uline{NB} Si $L_0$ n'est pas libre, alors $L_0^*$ est vide, et ce qui suit sans intérêt. On suppose donc par la suite $L_0$ \deleted{\unsure{dual|trivial}} \inserted{libre}, donc $L_1$ libre}

on a alors

\[
\left\{\begin{aligned}
 \sigma_0(\beta)^2&=\bigl(\begin{smallmatrix}-1&0\\0&-1\end{smallmatrix}\bigr)=u(-1)\overset{\mathrm{def}}{=}\unsure{\omega|w}
 \quad\text{(élément \uline{central} de $H$)}\\
 \sigma_0(\beta)(u(\lambda))&=u(\lambda^{-1})\quad\text{donc }v(\beta)\text{ normalise }T\\
 \sigma_0(\lambda\beta\deleted{\unsure{\lambda|*}})&=u(\lambda)\sigma_0(\beta)
 =\underset{\deleted{\text{\unsure{donc pour}}}}{\sigma_0(\beta)u(\lambda^{-1})},
\end{aligned}\right.\tag*{(20)}
\]

la dernière formule montre que, si on pose

\[
\deleted{\unsure{L_0}}N^-=v_0(L_0^*)=\deleted{H}\left\{\bigl(\begin{smallmatrix}a&b\\c&d\end{smallmatrix}\bigr)\in H\mid a=0,d=0\right\},\tag*{(21)}
\]

alors $N^-$ est une classe à gauche ou à droite mod $T$, et que

\[
L_0^*\xrightarrow{\sim}N^-\tag*{(22)}
\]

est bijection, et \={m} un isomorphisme de $\Lambda^*$-torseurs, quand on fait opérer $\Lambda^*$ sur $L_0^*$ par homothéties, et sur $N^-$ par multiplication \uline{à gauche} par $u(\lambda)$. On pourrait bien sûr introduire \deleted{aussi le composé} \inserted{symétriquement}

\[
\sigma_1:L_1\xrightarrow{\sim}N^-\qquad \sigma_1(\gamma)=\bigl(\begin{smallmatrix}0&-\gamma^{-1}\\\gamma&0\end{smallmatrix}\bigr)\tag*{\deleted{(23)} (23)}
\]

de sorte que l'on a

\[
\left\{\begin{aligned}
 \sigma_1(\gamma)&=\sigma_0(-\gamma^{-1})\\
 \sigma_1(\lambda\gamma)&=u(\lambda^{-1})\sigma_1(\gamma)=\sigma_1(\gamma)u(\lambda)
\end{aligned}\right.\tag*{(24)}
\]

donc $\sigma_1$ est un isomorphisme de \inserted{$\Lambda^*$-}torseurs, quand cette fois on fait opérer $\Lambda^*$ sur $N^-$ par multiplication à \uline{droite} par $u(\lambda)$.

% archive sheet number: 32
\scanpage{34}\mspage{3}

Notons la formule fondamentale

\[
\boxed{\sigma_0(\beta)(e_0(\beta))=\deleted{\unsure{\bigl(\begin{smallmatrix}1&0\\-\beta^{-1}&1\end{smallmatrix}\bigr)}=}\deleted{\unsure{u}}e_1(-\beta^{-1})}\tag*{(25)}
\]

% slanted
\margin{d'où plus généralement $\sigma_0(\beta)(e_0(\lambda\beta))=e_1(-\lambda\beta^{-1})$}

qui \deleted{\unsure{qui}} montre en particulier que

\[
\begin{aligned}
 \deleted{\unsure{tout|pour}}\forall\sigma\in N^-,\quad &\text{on a }\sigma(U_0)=\deleted{\unsure{u}}U_1\\
 &\text{d'où }\sigma(U_1)=U_0\qquad\deleted{\text{\unsure{pa}}}
\end{aligned}\tag*{(26)}
\]

(puisque $\sigma^2=\unsure{\omega|w}$ est central).

Notons que

\[
g=\bigl(\begin{smallmatrix}a&b\\c&d\end{smallmatrix}\bigr)\in H\quad\Longrightarrow\quad g^{-1}=\bigl(\begin{smallmatrix}d&-b\\-c&a\end{smallmatrix}\bigr)\tag*{(27)}
\]

et

\[
g e_0(\beta)g^{-1}=\bigl(\begin{smallmatrix}1-ac\beta&a^2\beta\\-c^2\beta&1+ac\beta\end{smallmatrix}\bigr)\tag*{(28)}
\]

donc on a $g(U_0)\subset U_1$ sss on a

\[
\begin{aligned}
 &a^2\beta=0,\ ac\beta=0\qquad\forall\,\beta\in L_0\\
 \text{i.e.}\qquad &a^2=0,\ ac=0\\
 \text{d'où}\qquad &a=a(\underbrace{ad-bc}_{1})=0
\end{aligned}
\]

donc

\[
\left\{\begin{aligned}
 \operatorname{Transp}_H(U_0,U_1)&=\left\{\bigl(\begin{smallmatrix}a&b\\c&d\end{smallmatrix}\bigr)\in H\mid a=0\right\},\ \text{de même}\\
 \operatorname{Transp}_H(U_1,U_0)&=\left\{\bigl(\begin{smallmatrix}a&b\\c&d\end{smallmatrix}\bigr)\in H\mid d=0\right\}.
\end{aligned}\right.\tag*{(29)}
\]

% slanted
\margin{il suffit même qu'un $g\in H$ soit tel que pour \uline{un} $\beta\in L_0^*$, on ait $g(\beta)\in U_1$, pour avoir $g\in\operatorname{Transp}_H(U_0\unsure{,|\to}U_1)$}

\deleted{\unsure{et}} \deleted{Je dis que} En fait, \unsure{le|les} \deleted{\unsure{on voit}} \inserted{calcul} \deleted{\unsure{on}} montre que dans ces cas là, le Transporteur large est strict, car si $a=0$, la formule (28) donne

\[
g e_0(\beta)g^{-1}=\bigl(\begin{smallmatrix}1&0\\-c^2\beta&1\end{smallmatrix}\bigr)\tag*{$(28')$}
\]

or $ad-bc=-bc=1$, donc $b$ et $c$ sont inversibles, donc $c^2$ inversible, donc $\operatorname{int}(g)$ induit un \uline{iso} $U_0\simeq U_1$. Des formules (29) on tire

% archive sheet number: 33
\scanpage{35}

\[
 N^-=\operatorname{Transp}(U_1,U_0)\cap\operatorname{Transp}(U_0,U_1)\tag*{(30)}
\]

Soient maintenant $\beta\in L_0$, $\gamma\in L_1$ avec

\[
\left\{\begin{aligned}
 &\beta\gamma=-1\quad\text{i.e. }\gamma=-\beta^{-1},\ \text{i.e. }\beta=-\gamma^{-1}\\
 &\text{i.e. }\ e_1(\gamma)=\underbrace{\sigma_0(\beta)(e_0(\beta))}_{e_1(-\beta^{-1})},
\end{aligned}\right.\tag*{(31)}
\]

je dis qu'on a alors

\[
 e_1(\gamma)e_0(\beta)e_1(\gamma)=\sigma_0(\beta)\qquad(\text{pour }\underline{\beta\gamma=-1})\tag*{(32)}
\]

formule qui s'écrit aussi, puisque $e_1(\gamma)=\sigma_0(\beta)e_0(\beta)\sigma_0(\beta)^{-1}$

\[
 \underbrace{\sigma_0(\beta)e_0(\beta)}\,\underbrace{\sigma_0(\beta)^{-1}e_0(\beta)}\,\underbrace{\sigma_0(\beta)e_0(\beta)}\,\sigma_0(\beta^{\deleted{\unsure{-1}}})^{-1}=\sigma_0(\beta)
 \tag*{\deleted{\unsure{(33)}}}
\]

soit encore, en écrivant $\sigma_0(\beta)^{-1}=\unsure{\omega|w}\,\sigma_0(\beta)$ ($\omega$ central, $\omega^2=1$)

\[
 (\sigma_0(\beta)e_0(\beta))^3\deleted{\unsure{\omega}}=\deleted{\unsure{\omega}}1
\]

soit encore

\[
 \rho(\beta)^3=1,\quad\text{où }\rho(\beta)=\sigma_0(\beta)e_0(\beta),\tag*{(33)}
\]

soit

\[
 \rho(\beta)=\bigl(\begin{smallmatrix}0&\beta\\-\beta^{-1}&-1\end{smallmatrix}\bigr)\deleted{\unsure{\text{est}}}
 \overset{\deleted{\text{soit }\rho(\beta}}{\underset{\substack{\text{si }L_0=L_1=\Lambda\\\text{donc }\beta\in\Lambda^*}}{=}}
 u(\beta)(\rho')\deleted{\unsure{u}}\quad\rho'=\rho(1)=\bigl(\begin{smallmatrix}0&1\\-1&-1\end{smallmatrix}\bigr)\tag*{(34)}
\]

la vérification est immédiate.

Supposons maintenant $\beta\in L_0$, $\gamma\in L_1$ quelconques, et calculons $e_1(\gamma)e_0(\beta)e_1(\gamma)$, on trouve

\[
 e_1(\gamma)e_0(\beta)(e_1(\gamma))=\bigl(\begin{smallmatrix}1+\beta\gamma&\beta\\\gamma(2+\beta\gamma)&1+\beta\gamma\end{smallmatrix}\bigr)\tag*{\unsure{(36)|(35)}}
\]

% slanted
\margin{faire le lien avec la formule pour $e_1(\deleted{\unsure{\gamma}}\beta)e_0(\deleted{\unsure{\beta}}\alpha')$ comme $e_0(x)u(t)e_1(y)$ pour $1+\deleted{\unsure{\alpha}}\beta\alpha'$ \unsure{invt.}}

qui prouve

\[
\begin{aligned}
 &\text{Si }\beta\in L_0,\gamma\in L_1,\text{ les conditions suivantes sont équivalentes}\\[-1pt]
 &\left\{\begin{aligned}
 \text{a)}\quad&\beta\gamma=-1\\
 \text{b)}\quad&e_1(\gamma)e_0(\beta)e_1(\gamma)\in\operatorname{Transp.}(U_0,U_1)\\
 \text{c)}\quad&e_1(\gamma)e_0(\beta)e_1(\gamma)\in\operatorname{Transp.}(U_1,U_0)\\
 \text{d)}\quad&e_1(\gamma)e_0(\beta)e_1(\gamma)\in N^-\qquad
 \text{e)}\quad e_1(\gamma)e_0(\beta)e_1(\gamma)=\sigma_0(\beta)
 \end{aligned}\right.
\end{aligned}\tag*{(37)}
\]

% archive sheet number: 34
\scanpage{36}\mspage{4}

En d'autres termes, en termes du groupe $H$, et des deux sous-groupes $U_0$, $U_1$ de $H$ et de leurs structures de $\Lambda$-modules inversibles $L_0$, $L_1$ (automatique si $\Lambda$ est un anneau absolu $\mathbb Z$ ou $\mathbb Z/n\mathbb Z$), on trouve

a) L'accouplement $L_0\otimes L_1\simeq\Lambda$, par la condition que $u_0\otimes u_1=-1$ sss

\[
\begin{aligned}
 &u_1u_0u_1\in\operatorname{Transp.}(U_0,U_1).\quad\text{ou encore}\\
 &u_1u_0u_1\in\operatorname{Transp}(U_0,U_1)\cap\operatorname{Transp}(U_1,U_0)\overset{\mathrm{def}}{=}N^-
\end{aligned}
\]

\margin{ou encore $u_0\in U_0^*$ et \deleted{$(u_1u_0u_1)($} $\operatorname{int}(u_1u_0u_1)(u_0)=u_1$ i.e. (38) ci-dessous}

b) L'application $\sigma_0:U_0^*\longrightarrow N^-$ par la condition $\sigma_0(u_0)=u_1u_0u_1$, où $u_1\in U_1$ est l'unique élément de $U_1$ satisfaisant à la condition que $u_1u_0u_1\in N^-$, et alors on a

\[
\text{c)}\qquad u_1=\sigma_0(u_0)(u_0)=\underbrace{u_1u_0u_1}u_0\underbrace{u_1^{-1}u_0^{-1}u_1^{-1}}
\]

i.e.

\[
 u_0u_1u_0u_1^{-1}u_0^{-1}u_1^{-1}=1\quad\text{i.e }\ u_0u_1u_0=u_1u_0u_1\tag*{(38)}
\]

formule qui est équivalente encore à $u_0\otimes u_1=-1$. En termes de $e_0,e_1$, (38) s'écrit aussi

\[
\begin{aligned}
 e_0(\beta)e_1(\gamma)e_0(\beta)=e_1(\gamma)e_0(\beta)e_1(\gamma)&\qquad\text{sss }\beta\gamma=-1\\
 &(\beta\in L_0^*,\ \gamma\in L_1).
\end{aligned}\tag*{(39)}
\]

La relation (38) \deleted{est} entre $U_0$ et $U_1$ est le graphe d'une bijection entre $U_0^*=e_0(L_0^*)$ et $U_1^*=e_1(L_1^*)$, qu'on peut noter

\[
u_0\longmapsto u_0^\vee:U_0^*\longrightarrow U_1^*
\]

et satisfaisant

\[
(\deleted{\unsure{\lambda}}u_0^\lambda)^\vee=(u_0^\vee)^{\lambda^{-1}},\qquad u_0\in U_0^*,\ \lambda\in\Lambda^*\tag*{(40)}
\]

en notant exponentiellement l'opération \inserted{externe} de $\Lambda$ sur $U_0,U_1$.

% archive sheet number: 35
\scanpage{37}

(alors que pour $g\in\operatorname{Transp}_H(U_0,U_1)$, \deleted{\unsure{le|l'i}}la bijection

\[
\operatorname{int}(g)|U_0:\ U_0\simeq U_1
\]

respecte les structures de $\Lambda$-module sur $U_0$ et $U_1$ cf $(28')$).

\deleted{On notera}

\[
\deleted{N'=T\sqcup N^-\subset H}\tag*{\deleted{(41)}}
\]

\deleted{\unsure{c'}est un sous-groupe de}

% archive sheet number: 36
\scanpage{38}

Rapport anharmonique

\unsure{avri|avril} 1986

% archive sheet number: 37
\scanpage{39}\mspage{1}

Avril 86

À propos rapport anharmonique.

\textcircled{a)} Soi\deleted{\unsure{ent}}t $V\xrightarrow{\varphi_i}V_i$ ($i\in I$, $\operatorname{card} I=3$) des flèches de $\unsure{=|\overline{m}}$ source dans une cat. abélienne \inserted{$C$}, telles que $\forall J\subset I$, $\operatorname{card} J=2$, la flèche correspondante \deleted{soit}

\[
 V\longrightarrow V_J\overset{\mathrm{def}}{=}\prod_{i\in J}V_i
\]

soit iso. \deleted{\unsure{Cons|Considérons}} Alors on a suite exacte

\diagram{\[
\begin{tikzcd}[ampersand replacement=\&,column sep=small,row sep=small]
 0\arrow[r]\& V\arrow[r,"\varphi","{\varphi=(\varphi_i)}"']\& \prod V_i\arrow[r]\& D\arrow[r]\arrow[d,equal]\&0\\
 \&\&\&\prod V_i/\operatorname{Im}V\&
\end{tikzcd}\tag*{(1)}
\]}

et les \unsure{trois|hom} composés

\diagram{\begin{tikzcd}[ampersand replacement=\&]
V_i\arrow[r,"\mathrm{inj}_i"]\arrow[rr,bend right=20,"u_i"']\&\prod V_i\arrow[r]\&D
\end{tikzcd}}

sont iso. On a donc un iso. de suites exactes entre (1) et

\diagram{\begin{tikzcd}[ampersand replacement=\&,column sep=small]
0\arrow[r,dash]\&\deleted{\unsure{V}}K(D,I)\arrow[r]\&D^I\arrow[r,"{\varepsilon_{D,I}}","\mathrm{somme}"']\&D\arrow[r]\&0
\end{tikzcd}}

\[
 \text{où}\quad K(D,I)\overset{\mathrm{def}}{=}\operatorname{Ker}\varepsilon_{D,\unsure{E|I}}
 =\left\{(x_i)_{i\in I}\in D^I\mid\sum x_i=0\right\}
\]

Ainsi, la catégorie des systèmes $(V,(V_i),(\varphi_i))$ équivaut à celle des paires $(D,I)$, i.e. est équivalente à $\mathrm{Ens}_3\times C$.

\textcircled{b} \uline{NB} Il y a une variante non commutative ensembliste, \deleted{avec \unsure{détermination|déterminations|d'éventuelles lois|dimensions} de} \inserted{où $V,V_i$ sont des ens sans} plus. \deleted{\unsure{On}} On doit pouvoir montrer

% archive sheet number: 38
\scanpage{40}\mspage{2}

(utilisant une orientation de $I$) qu'il existe sur chaque $V_i$ une structure de bitorseur \inserted{\unsure{ens|sur|sous}} sous une paire de groupes $G_{\underset{\substack{\Vert\\\rho^{-1}(i)}}{j}},G_{\underset{\substack{\Vert\\\rho(i)}}{k}}$ (les $G_j$ étant définis \deleted{\unsure{p}} à isom unique près)

\diagram{\begin{tikzcd}[ampersand replacement=\&,column sep=small,row sep=small]
 \& k\;\;V_k\arrow[dl,dash,"G_j"']\arrow[dr,dash,"G_i"]\&\\
 V_i\;\;i\arrow[rr,dash,"G_k"']\&\&j\;\;V_j
\end{tikzcd}}

et un isom (choisissant un $i\in I$ pour la commodité de l'écriture

\[
V_i\mathbin{\underset{G_k\deleted{\unsure{i}}}{\wedge}}V_j
 \mathbin{\underset{G_i}{\wedge}}V_k\simeq\mathbb{1}_{G_i}
 \qquad\text{isom de bitorseurs}
\]

i.e. un élément central \inserted{$\varepsilon_i$} de $V_i\mathbin{\underset{G_k}{\wedge}}V_j\mathbin{\underset{G_i}{\wedge}}V_k$ de telle façon que $V$ s'identifie à l'ens des $(x_i,x_j,x_k)$ tels que

\[
 x_i\wedge x_j\wedge x_k=\varepsilon
\]

Si les $G_i$ sont \deleted{commutatifs, ils sont} \inserted{\unsure{can}$^{\mathrm t}$. ext. \unsure{isop.}}, donc \unsure{si|où} l'un est comm, les autres le sont et on a un syst. transitif d'isomorphismes ; On est ramené à la situation précédente, dès qu'on se donne une ``origine'' $o\in V\ldots$

% archive sheet number: 39
\scanpage{41}\mspage{3}

\deleted{\unsure{c}} \uline{Remarque} Le groupe des automorphismes de la situation donnée se dévisse en $\unsure{\mathfrak S|\mathscr S}_I$ ($\simeq\mathfrak S_3$) \deleted{$\operatorname{Aut}G_j\times_{\unsure{\operatorname{Aut}\operatorname{Ext}G_i}}\operatorname{Aut}G_i$} et

\[
\underset{\simeq\Gamma\ \text{si les }G_i\ \text{\unsure{comm.}}}{\underbrace{\prod_{i\in I}{}_{\Gamma}\operatorname{Aut}(G_{i})}}\,,
\quad\text{où }\Gamma^{\deleted{\unsure{\varepsilon}}}=
 \underbrace{\operatorname{Aut}\operatorname{Ext}(G_i)}_{\substack{\text{indépendant de }i\\\text{mod syst. tran.}\\\text{d'iso}}}
\]

$\prod G_i$, opérant sur $(V_i)$ via

\[
g_jx_i g_k^{-1},\quad g_kx_jg_i^{-1},\quad g_ix_kg_j^{-1}
\]

% slanted
\margin{\unsure{oubli !|voir !|restriction|produit} au vrai \unsure{compte|conjugués}, il faut prendre \deleted{$\unsure{G^{I}/Z}\ \unsure{Z^I/\mathrm{diag}|Z^I/\mathrm{diagonale}}$} où $Z$ est le ``centre commun'' des $G_i$}

\margin{Il est plus joli de dévisser en $\mathfrak S_I$, \deleted{Aut}$\operatorname{Ext}(G_{\unsure{i|*}})$, et $G^I/Z$ (envoyé diagonalmt en $Z^I\subset G^I$}

\textcircled{c} Soit maint$^{\mathrm{nt}}$ $I$ \unsure{un|ens} à \uline{quatre} éléments, et \deleted{\unsure{d'un}} \unsure{se|on|on se} \unsure{proposons|propose} \unsure{nous|encore} d'étudier les

\[
 V\xrightarrow{\varphi_i}V_i\qquad i\in I
\]

dans $C$, tels que $\forall J\in I$, $\operatorname{card} J=3$, on ait

\[
V\xrightarrow{\varphi_J}V_J\qquad\text{iso.}
\]

Pour toute \deleted{subdivision} numérotation

\[
 I=\{i_1,i_2,i_3,i_4\}
\]

on va définir une flèche

\[
 u_{(i_1,i_2),(i_3,i_4)}:V_{i_1}\xrightarrow{\sim}V_{i_2}
 \tag*{(1)}
\]

\margin{\unsure{24|Il} \unsure{quatrts|quantités|quadruplets|par}}

C'est l'iso can. défini en a), quand on fait abstraction de $\varphi_{i_4}$, \deleted{\unsure{Il|donc}}

% archive sheet number: 40
\scanpage{42}\mspage{4}

\[
u_{(i_2,i_1)(\unsure{i_3},i_4)}
 =u_{(i_1,i_2)(i_3,i_4)}^{-1}\tag*{(2)}
\]

\margin{12 relations reste à fixer 12 quantités}

\[
\begin{aligned}
 &\deleted{\unsure{u}_{\unsure{i_1i_2}}}\\
 &\deleted{u_{(i_2\,i_1)(i_3,\unsure{i})}}
\end{aligned}
\]

\[
\deleted{u_{i_2\,i_1\,i_3\,i_4}}\,
\unsure{u_{(i_1\,i_3)(i_2\,i_4)}}\,
u_{(i_3,i_2)\unsure{(i_1\,i_4)|(i_1\,i_2)}}\,u_{(i_2\,i_1)(i_3\,i_4)}=\mathrm{id}\tag*{(3)}
\]

$V$ est connu en tant que sous-objet de $\prod V_i$, quand on connait le système des $u$, car c'est formé des $(x_i)_{i\in I}$ tels que pour tout $i\in I$ \deleted{\unsure{chaque|donné}} les $x_j$ ($j\in I\setminus\{i\}$) soient ``de somme nulle'' pour le système transitif d'iso entre ces $V_j$ ($j\in J$) défini par l'une ou l'autre (ça revient au $=$) des deux perm. circulaires de $J$.

Mais il faut des conditions pour que le $W$ ainsi défini (qui a priori sera tel que $W\longrightarrow V_i\times V_j=V_J$ \deleted{soit} soit injectif) soit ``assez gros'' pour que ce soit bijectif. On va écrire des conditions pour ceci.

% archive sheet number: 41
\scanpage{43}\mspage{5}

On va définir, pour un système d'iso $u_{ij\unsure{k|h}\unsure{l|\ell}}$, un système de ``rapports anharmoniques''

\[
\rho_{ijkl}:\operatorname{Aut}V_{\unsure{0|i}}\qquad(\{i,j,k,l\}=I\}\tag*{\unsure{$(4_0)$|(4)}}
\]

On a

\[
\begin{aligned}
u_{ijkl}&: V_i\xrightarrow{\sim}V_j\\
u_{ijlk}&: V_i\xrightarrow{\sim}V_j
\end{aligned}
\]

\margin{\unsure{$\times$}\quad\unsure{$'$}\quad .\quad ?\quad .}

d'où

\[
 \rho_{ijkl}=u_{ijlk}^{-1}u_{ijkl}\deleted{\unsure{e}},\tag*{(4)}
\]

\margin{24 quantités}

Donc on aura

\[
 \rho_{ijlk}=(\rho_{ijkl})^{-1}\tag*{(5)}
\]

\margin{\deleted{12} relations \unsure{donc|dont|restent} 12 quantités correspondant aux arêtes orientées de $I$}

On aura d'ailleurs

\[
\underset{\substack{\Vert\\u_{jilk}^{-1}u_{jikl}}}{\rho_{jikl}}=
\deleted{\unsure{u_{ijkl}\rho_{ijkl}}}
\left\{\begin{aligned}
 u_{ij\unsure{kl}}(\rho_{ijkl}^{-1})\deleted{\unsure{u_{ijkl}^{-1}}}\\
 u_{ijlk}(\rho_{ijkl}^{-1})
\end{aligned}\right.\tag*{(6)}
\]

\margin{NB On pose $\beta(\alpha)=\beta\alpha\beta^{-1}$}

[ Ce qui donne

\[
\begin{aligned}
\rho_{j\unsure{,}i,l,k}&=u_{ijkl}(\rho_{ijkl})\\
 &=u_{ijlk}(\rho_{ijkl})
\end{aligned}
\]

\margin{Invariance par double permutation}

% archive sheet number: 42
\scanpage{44}\mspage{6}

Donc il reste essentiellement \deleted{\unsure{3|8}}6 quantités, i.e. \deleted{pour} il suffit, pour \deleted{\unsure{le|tt}} \deleted{\unsure{partitio|partition}} \deleted{subdivision $(2,2)$ \unsure{ti} de $I$ ; dis. $I'\sqcup I''$,} partie à deux éléments $\underset{\substack{\Vert\\\{i'_1,i'_2\}}}{I'}\subset \underset{\substack{\Vert\\\{i''_1,i''_2\}}}{I^{\unsure{2|*|\prime\prime}}}$, de de déterminer l'une des quatre quantités

\[
 \rho_{\underset{\tikz{\draw[<->] (0,0) to[bend right=35] (.36,0);}}{i'_1\,i'_2}\;
 \underset{\tikz{\draw[<->] (0,0) to[bend right=35] (.36,0);}}{i''_1\,i''_2}}.
 \qquad\deleted{\unsure{Il}}
\]

$\star$ Nous allons \deleted{\unsure{don}} enfin donner une formule pour relier

\[
\rho_{ij\unsure{k|h}\unsure{l|\ell}}\quad\text{et}\quad\rho_{i\deleted{\unsure{j}}kjl}
\]

(ce qui permettra donc de passer d'un repère à un repère quelconque)

\[
 \rho_{ijkl}+\rho_{ikjl}=1
\]

i.e

\[
\boxed{\rho_{i\deleted{\unsure{jk}}kjl}=1-\rho_{i\deleted{\unsure{kj}}jkl}}\tag*{(7)}
\]

\margin{[ Faire vérification}

\uline{Prop} Soit un système de $u_{ijkl}:V_i\longrightarrow V_j$, satisfaisant les relations (2) (3) ; d'où (4) des $\rho_{ijkl}\in\operatorname{Aut}(V_i)$, \unsure{)}, satisfaisant (5)(6). Pour que ça corresponde à une \unsure{quadrité|quadrite|quadralité} \deleted{\unsure{$V$}} $(V,\deleted{\unsure{\varphi}}(V_i),(\varphi_i))$ \deleted{\unsure{$\subset$}} $V\subset\prod V_i$, il

% archive sheet number: 43
\scanpage{45}
\mspage{7}
% archive sheet number: 44

faut et il suffit qu'on ait aussi les 24 relations (7). Quand \inserted{les} $\deleted{\unsure{\operatorname{Aut}}}\operatorname{End}(V_i)$ sont commutatifs (rappelons qu'ils sont iso entre eux) alors \deleted{les \unsure{isom|isoms}} \inserted{les seules} \deleted{$\operatorname{Aut} V_i$} relations (2), (3) impliquent que l'on a un syst. trans. d'iso. entre les \inserted{anneaux} $\deleted{\unsure{\operatorname{Aut}}}\operatorname{End}(V_i)$, soit $\deleted{\unsure{A}}\unsure{k|h}$ \deleted{\unsure{est}}\unsure{cet} anneau. Alors les relations 5, 6, 7 s'interprètent comme des relations dans \deleted{\unsure{E|Ens}} $k$

\[
\left\{\begin{aligned}
\unsure{\rho|f}_{\deleted{\unsure{ij}}j i k\unsure{l|\ell}}&=\rho_{ijlk}=\rho_{ijkl}^{-1} \qquad (\Rightarrow \rho_{\deleted{\unsure{ij}}jilk}=\rho_{ijkl})\\
\deleted{\unsure{\rho}}\rho_{ikjl}&=1-\rho_{ijkl}
\end{aligned}\right. \tag*{(8)}
\]

En d'autres termes, faisant opérer $\unsure{\mathfrak S|\mathscr S}_3$ sur l'ens \inserted{$U_{\deleted{\unsure{k}}}(k)$} des $\lambda\in k$ tels que $\lambda$ et $1-\lambda$ sont inversibles, par

\[
\begin{aligned}
\sigma_{12}(\lambda)&=\sigma_{34}(\lambda)=\lambda^{-1}\\
\sigma_{23}(\lambda)&=1-\lambda
\end{aligned}
\]

\margin{donc $\mathfrak S_4$ opère par l'intermédiaire du quotient $\mathfrak S_3$ de $\mathfrak S_4$}

\deleted{\unsure{avec} $\unsure{\rho}$} l'application \deleted{$\rho_{ijkl}$}

\[
\underset{\text{torseur sous }\mathfrak S_4}{\operatorname{Rep}(I)}\longrightarrow \deleted{\unsure{k}}U(k)
\qquad (i,j,k,l)\unsure{\text{---}|\longrightarrow}\rho_{ijkl}
\]

est comp. aux $\mathfrak S_4$
